% Kelvin. This is also the control for a11.
\subsection{\expt{a14}: Maintenance before benchmarking}
Experiment a14 runs \texttt{VACUUM (ANALYZE)} on each table after loading, followed by
\texttt{CHECKPOINT} on all five nodes. The motivation was to reduce interference
from dead tuples and pending writes left by data generation. It also provides
the maintenance control for a11.

Job 967194 showed no consistent overall TPS or transaction-latency benefit.
At eight clients, T9's median decreased but T8's increased. One final
serialization failure occurred at 32 clients. The experiment did not justify
including this procedure in the final A design; its principal value was as
a control for fillfactor.
\evidence{\expt{a14}/job 967194: \code{hot_updates.csv},
\code{transaction_metrics.csv} and \code{summary.csv}.}